% =============================================================
%  Module Description — Calculus II (MAT231-F27)
%  Shown on https://wosm.academy as the course page for MAT231-F27.
%  Compile with: pdflatex mat231-module-description.tex
%  Write "TODO(owner): ..." where information is still missing;
%  the website highlights it.
% =============================================================
\documentclass[11pt,a4paper]{article}

\usepackage[utf8]{inputenc}
\usepackage[T1]{fontenc}
\usepackage{lmodern}
\usepackage[margin=2.2cm]{geometry}
\usepackage{amssymb}
\usepackage{booktabs}
\usepackage{tabularx}
\usepackage{array}
\usepackage{enumitem}
\usepackage{xcolor}
\usepackage{titlesec}
\usepackage{fancyhdr}
\usepackage[hidelinks]{hyperref}

% ---------- Module details (edit these) ----------------------
% \ModuleCode must equal the course's `code` in src/content/courses/.
\newcommand{\ModuleTitle}{Calculus II}
\newcommand{\ModuleCode}{MAT231-F27}
\newcommand{\Programme}{Foundation programme}
\newcommand{\Level}{Semester 3}
\newcommand{\Semester}{Fall 2027}
\newcommand{\Credits}{7 ECTS}
\newcommand{\ModuleLeader}{Dr. Ali Khudiyev}
\newcommand{\Email}{ali.khudiyev@wosm.academy}
\newcommand{\OfficeHours}{Monday--Friday, 10:00--17:00 (Baku time, UTC+4)}
\newcommand{\Prerequisites}{MAT102-S27 Linear Algebra; MAT103-S27 Calculus I}

% ---------- Styling ------------------------------------------
\definecolor{accent}{RGB}{0,70,130}
\titleformat{\section}{\large\bfseries\color{accent}}{\thesection.}{0.5em}{}
\titlespacing*{\section}{0pt}{1.2em}{0.5em}
\setlist{itemsep=2pt, topsep=4pt}
\setlength{\parindent}{0pt}
\setlength{\parskip}{4pt}

\pagestyle{fancy}
\fancyhf{}
\lhead{\small \ModuleCode\ -- \ModuleTitle}
\rhead{\small Module Description}
\cfoot{\small \thepage}
\renewcommand{\headrulewidth}{0.4pt}

% =============================================================
\begin{document}

\begin{center}
  {\LARGE\bfseries\color{accent} \ModuleTitle}\\[4pt]
  {\large Module Description}
\end{center}
\vspace{6pt}

% ---------- Summary table (not shown on the website: the course page
% shows these details from the macros above) --------------------------
\begin{tabularx}{\textwidth}{>{\bfseries}l X >{\bfseries}l X}
\toprule
Module code   & \ModuleCode    & Credits      & \Credits \\
Programme     & \Programme     & Level        & \Level \\
Semester      & \Semester      & Prerequisites & \Prerequisites \\
Module leader & \ModuleLeader  & Email        & \href{mailto:\Email}{\Email} \\
Office hours  & \multicolumn{3}{l}{\OfficeHours} \\
\bottomrule
\end{tabularx}

% ---------- Sections (shown on the website) ------------------
\section{Module Overview}
Calculus II continues Calculus I in two directions. First it completes the
one-variable theory with infinite series, power series and Taylor series. Then it
extends differentiation and integration to functions of several variables, using
the vectors and matrices of Linear Algebra.

The module covers:
\begin{itemize}
  \item Infinite series, convergence tests, power series and Taylor series
        (first 4--5 weeks)
  \item Differential calculus of several variables: partial derivatives, gradients,
        the Jacobian and Hessian, and optimisation with and without constraints
  \item Multiple integrals and change of variables
  \item An introduction to ordinary differential equations
\end{itemize}

These tools are used directly later in the programme: gradients and Hessians in
optimisation and machine learning, multiple integrals in Probability Theory and
Statistics (joint densities), and series and generating functions in the analysis
of algorithms.

\section{Aims}
\begin{itemize}
  \item Complete the theory of one-variable calculus with series and uniform convergence.
  \item Extend differential and integral calculus to functions of several variables.
  \item Give students the tools for continuous optimisation, including gradient-based methods.
  \item Prepare students for probability, statistics, machine learning and the analysis of algorithms.
\end{itemize}

\section{Learning Outcomes}
On successful completion of this module, students will be able to:
\begin{enumerate}[label=LO\arabic*.]
  \item Decide the convergence of series with the comparison, ratio, root, integral and alternating-series tests, and find the radius of convergence of power series.
  \item Expand functions as Taylor series, bound the remainder, and use series and generating functions to solve counting and recurrence problems.
  \item Compute partial derivatives, gradients, directional derivatives, Jacobians and Hessians, and apply the multivariable chain rule.
  \item Find and classify critical points, solve constrained problems with Lagrange multipliers, and explain and apply gradient descent.
  \item Evaluate double and triple integrals, change the order of integration, and change variables using polar, cylindrical and spherical coordinates and the Jacobian determinant.
  \item Solve first-order and linear second-order ordinary differential equations and simple linear systems using eigenvalues.
\end{enumerate}

\section{Syllabus}
\begin{tabularx}{\textwidth}{c X}
\toprule
\textbf{Unit} & \textbf{Topic} \\
\midrule
1  & Infinite series; geometric and telescoping series \\
2  & Convergence tests: comparison, integral, ratio, root \\
3  & Absolute and conditional convergence; alternating series; rearrangements \\
4  & Sequences and series of functions; pointwise and uniform convergence \\
5  & Power series; radius of convergence; term-by-term differentiation and integration \\
6  & Taylor and Maclaurin series; remainder estimates \\
7  & Generating functions; solving recurrences \\
8  & Stirling's approximation \\
9  & $\mathbb{R}^n$: norms, open and closed sets, limits and continuity \\
10 & Partial derivatives; differentiability; the total derivative \\
11 & Gradient and directional derivatives; level sets \\
12 & The multivariable chain rule; the Jacobian matrix \\
13 & Higher derivatives; the Hessian; multivariable Taylor expansion \\
14 & Unconstrained optimisation: critical points, second-derivative test, convexity \\
15 & Gradient descent and its convergence for convex quadratics \\
16 & Constrained optimisation: Lagrange multipliers \\
17 & Implicit and inverse function theorems (statement and use) \\
18 & Double integrals over rectangles and general regions; Fubini's theorem \\
19 & Triple integrals \\
20 & Change of variables; polar, cylindrical and spherical coordinates \\
21 & The Gaussian integral; improper multiple integrals \\
22 & Introduction to vector fields and line integrals \\
23 & First-order ODEs: separable and linear equations \\
24 & Second-order linear ODEs with constant coefficients \\
25 & Systems of linear ODEs via eigenvalues \\
\bottomrule
\end{tabularx}

\section{Teaching and Learning Methods}
The module is self-paced. All materials are provided on Moodle: lecture notes and
recorded lectures for each unit, tutorial problem sheets with worked solutions,
short quizzes, and the briefs for the two projects. Problem sheets are for practice
and are not marked. Questions are discussed in the module channel on Discord and in
office hours. Independent study includes reading, problem sheets and project work.

\begin{tabularx}{0.6\textwidth}{X r}
\toprule
\textbf{Activity} & \textbf{Hours} \\
\midrule
Lectures            & 50 \\
Labs / Tutorials    & 30 \\
Independent study   & 130 \\
\midrule
\textbf{Total}      & \textbf{210} \\
\bottomrule
\end{tabularx}

\section{Assessment}
Knowledge is assessed with weekly unit quizzes, two projects and a final
exam. Because much of the module is theory, the exam carries slightly more weight.
\begin{itemize}
  \item \textbf{Project 1 -- Series and generating functions:} use generating
        functions to solve recurrences from counting problems and algorithms, and
        check the results numerically.
  \item \textbf{Project 2 -- Optimisation in several variables:} analyse a function
        with gradients and Hessians, solve a constrained version with Lagrange
        multipliers, and implement gradient descent, studying how the step size
        affects convergence.
\end{itemize}
The final exam tests computation (series, derivatives and integrals in several
variables), the choice of method, and short proofs about convergence.

\begin{tabularx}{\textwidth}{X c c c}
\toprule
\textbf{Component} & \textbf{Weight} & \textbf{Due} & \textbf{LOs} \\
\midrule
Unit quizzes (best 10 of 12) & 20\% & Weekly & LO1--LO3, LO5 \\
Project 1: Series and generating functions & 15\% & Week 6 & LO1, LO2 \\
Project 2: Optimisation in several variables & 20\% & Week 12 & LO3, LO4 \\
Final exam (3 hours, written) & 45\% & End of term & LO1--LO6 \\
\midrule
\textbf{Total}      & \textbf{100\%} & & \\
\bottomrule
\end{tabularx}

\section{Reading List}
\textbf{Essential}
\begin{itemize}
  \item G.\ Strang, E.\ Herman: \textit{Calculus, Volumes 2 and 3}, OpenStax, 2016. Free online at \url{https://openstax.org/subjects/math}.
  \item J.\ E.\ Marsden, A.\ Tromba: \textit{Vector Calculus}, 6th ed., W.\ H.\ Freeman, 2012.
\end{itemize}
\textbf{Recommended}
\begin{itemize}
  \item T.\ M.\ Apostol: \textit{Calculus, Volume 2}, 2nd ed., Wiley, 1969.
  \item J.\ H.\ Hubbard, B.\ B.\ Hubbard: \textit{Vector Calculus, Linear Algebra, and Differential Forms}, 5th ed., Matrix Editions, 2015.
  \item H.\ S.\ Wilf: \textit{generatingfunctionology}, 3rd ed., A K Peters, 2006. Free online at \url{https://www2.math.upenn.edu/~wilf/DownldGF.html}.
  \item S.\ Boyd, L.\ Vandenberghe: \textit{Convex Optimization}, Cambridge University Press, 2004 (Chapter 9).
\end{itemize}

\section{Materials and Licence}
Course materials are published under the
\href{https://creativecommons.org/licenses/by-sa/4.0/}{Creative Commons Attribution-ShareAlike 4.0}
licence unless stated otherwise. See \href{https://wosm.academy/licence/}{wosm.academy/licence} for details.

\section{Policies}
\begin{itemize}
  \item \textbf{Attendance:} Attendance is not mandatory. The module is self-paced, so you can study the materials when it suits you.
  \item \textbf{Late submissions:} Work submitted after the deadline is penalised.
  \item \textbf{Academic integrity:} All submitted work must be your own.
        See the \href{https://wosm.academy/code-of-conduct/}{Code of Conduct}.
\end{itemize}

\end{document}
